%% ----------------------------------------------------------------
%% Week-10 problem sheet (lectures: entropy, information, wasserstein).
%% Every number in the answers was checked in numpy/PyTorch (2026-10-30).
%% ----------------------------------------------------------------
\documentclass{sotonExamBoxes}    % use option [answers] to include them
%% ----------------------------------------------------------------
\usepackage{graphicx}
\usepackage{alltt}
\usepackage{bm}
\usepackage{amssymb}
\usepackage{amsfonts}
\usepackage{amsmath}
\thinmuskip=1mu
%% macros.tex -- shared macro definitions for the problem sheets.
%%
%% Every sheet inputs this instead of carrying its own copy. Input it after
%% the \usepackage lines: the definitions below use bm, amsmath and amssymb.
%%
%% Six macros were defined inconsistently across the sheets when this file
%% was made; the choices are marked CONFLICT below.




\usepackage{stmaryrd}

\renewcommand{\ttdefault}{pcr}

\newcommand{\tr}{\textsf{T}}

\newcommand{\e}[1]{{\rm e}^{#1}}

\newcommand{\E}[1]{{\rm e}^{{\displaystyle #1}}}

\newcommand{\logg}[1]{\log\left(#1\right)}

\newcommand{\Prob}[1]{\mathbb{P}\left(#1\right)}

\newcommand{\Step}{\mathop{\mathrm{O}\hspace{-15pt}\rule[8.5pt]{11pt}{1.2pt}\hspace{2pt}}}

\newcommand\diag{\mathop{\mathrm{diag}}}

\newcommand{\grad}{\bm{\nabla}}

\newcommand{\dd}{\mathrm{d}}

\newcommand{\class}{\mathcal{C}}

\newcommand{\data}{\mathcal{D}}

%% CONFLICT: gaussianProcess.tex had {U}{eur}{b}{n} (Euler bold);
%% every other sheet had this. Matrices in that one sheet now look
%% different -- see the report.

\DeclareMathAlphabet{\mat}{OT1}{cmss}{bx}{n}

%% The identity needs to be distinguishable. In cmss a bold I is a plain
%% vertical bar with no crosspieces, so \mat{I} reads as a rule rather than a
%% letter. Euler's bold I keeps its serifs, so the identity alone is set in
%% Euler while everything else stays cmss.
\DeclareMathAlphabet{\matEuler}{U}{eur}{b}{n}
\newcommand{\Id}{\matEuler{I}}

%% CONFLICT and BUG FIX: 16 sheets defined this using \qpartial,
%% which is defined nowhere, so \pd was broken in all of them.
%% Only problem2.tex had the working \partial form, kept here.

\newcommand{\pd}[2]{\ensuremath{\frac{\partial #1}{\partial #2}}\xspace}

\newcommand{\M}[1]{\ensuremath{\boldsymbol{#1}}\xspace}

\newcommand{\bst}{\rule{0mm}{5mm}}

%% CONFLICT: some sheets omitted the \! . Kept the tighter spacing,
%% which is what every sheet that actually uses \av had.

\newcommand{\av}[2][\,]{\mathbb{E}_{#1}\!\left[ {\strut #2} \right]}

\newcommand{\len}[1]{\| #1 \|}

\newcommand{\inner}[2]{\left\langle #1, #2\right\rangle}

\newcommand{\squeeze}{\setlength{\itemsep}{-2pt}}

%% CONFLICT: two spacings, identical output in maths mode.

\newcommand{\bra}[1]{\left( #1 \right)}

\newcommand{\pred}[1]{\left\llbracket { \small #1} \right\rrbracket}

\newcommand{\mln}{\hat{f}_n(\bm{x})}


\newcommand{\Tr}{\mathop{\mathrm{tr}}\,}

%% NOT SHARED: gaussianProcess.tex redefines \det to take an argument
%% (\det{\mat{V}}), while calculus.tex and vectorSpaces.tex use the standard
%% no-argument operator (\det(\mat{X})). The two cannot coexist, so that
%% redefinition stays local to gaussianProcess.tex.

%% NOT SHARED: \ml and \mm mean different things in different sheets.
%% amgm and convexity use a no-argument \ml and \mm = \hat{m}_c(x);
%% ensembleLearning uses a one-argument \ml and \mm = \hat{m}(x).
%% Both forms are in use, so each sheet keeps its own.

\graphicspath{{figures/}}
\usepackage{parskip}
\AddToHook{env/samepage/begin}{%
\AddToHook{para/before}{\nopagebreak}%
\AddToHook{para/after}{\nopagebreak}%
}
\newcommand{\KL}[2]{\mathrm{KL}\bra{#1\,\|\,#2}}
\begin{document}
\unitTitle{Entropy, Information and Optimal Transport Problem Sheet}
\authors{Adam Pr\"ugel-Bennett}
\unitcode{AICE3001/AICE6002}
\semester{2}
\year{2022}
\duration{45}{0}

\maketitle

This sheet covers the lectures on entropy, information theory and
optimal transport.  Question~1 practises entropy and code lengths,
Question~2 mutual information, the KL divergence and cross-entropy, and
Question~3 the maximum entropy method.  Each ends with a short Python
part (\texttt{import numpy as np} is assumed).  Question~4 is harder and
optional: you write a small optimal transport solver in PyTorch and use
it to move points by gradient descent.  Logarithms are base 2 (bits)
unless we say otherwise.

% ***************************** question 1 **************************

\begin{question}{20}
  \begin{qparts}

    \qpart[4]{Find the entropy of (i) a fair coin, (ii) a coin with
      $\Prob{\text{head}} = 0.9$, and (iii) a random variable $X$ taking
      the values $a$, $b$, $c$, $d$ with probabilities $\frac12$,
      $\frac14$, $\frac18$, $\frac18$.
      \begin{answer}
        (i) $H = -2\times\frac12\log_2\frac12 = 1$ bit.
        (ii) $H = -0.9\log_2 0.9 - 0.1\log_2 0.1 \approx 0.137 + 0.332
        = 0.469$ bits: the outcome is much less uncertain.
        (iii) $H = \frac12\times1 + \frac14\times2 + 2\times\frac18\times3 =
        1.75$ bits.
      \end{answer}}{\lines{5}}

    \qpart[4]{We code the outcomes of $X$ in part (a)(iii) as $a\to0$,
      $b\to10$, $c\to110$, $d\to111$.  Why can a string of these code
      words always be decoded?  Decode \texttt{0110100111}.  Find the
      expected code length and compare it with the entropy.
      \begin{answer}
        It is a prefix code: no code word is the start of another, so
        reading from the left we always know where a code word ends.
        \texttt{0|110|10|0|111} is $a\,c\,b\,a\,d$.  The expected length
        is $\frac12\times1 + \frac14\times2 + \frac18\times3 + \frac18\times3
        = 1.75$ bits, exactly the entropy, because each code length is
        $-\log_2$ of its probability.
      \end{answer}}{\lines{5}}

    \qpart[4]{Instead we use two bits for every outcome
      (\texttt{00}, \texttt{01}, \texttt{10}, \texttt{11}).  How many
      extra bits per outcome does this cost?  What is the largest entropy
      any distribution on four outcomes can have, and which distribution
      has it?
      \begin{answer}
        The expected length is 2 bits, so $0.25$ extra bits per outcome.
        (This is the KL divergence between $X$'s distribution and the
        uniform distribution that the code is optimal for.)  Entropy is
        concave and, by Jensen's inequality, at most $\log_2 4 = 2$ bits,
        reached only by the uniform distribution: for uniform outcomes the
        two-bit code is optimal.
      \end{answer}}{\lines{5}}

    \qpart[3]{$X_1$ and $X_2$ are independent copies of $X$.  What is
      the entropy of the pair $(X_1, X_2)$?
      \begin{answer}
        For independent variables entropies add:
        $H = 1.75 + 1.75 = 3.5$ bits.
      \end{answer}}{\lines{3}}

    \qpart[5]{Write a Python function \texttt{entropy(p)} that returns the
      entropy in bits of a probability vector (treat $0\log0$ as 0).
      Check part (a), then draw $10\,000$ random distributions on four
      outcomes with \texttt{rng.dirichlet(np.ones(4))} and find the
      largest entropy.
      \begin{answer}
\begin{alltt}
def entropy(p):\\
\hspace*{2em}p = np.asarray(p, dtype=float)\\
\hspace*{2em}p = p[p > 0]\hspace*{2em}\# 0 log 0 = 0\\
\hspace*{2em}return -(p * np.log2(p)).sum()\\
print(entropy([0.5, 0.5]), entropy([0.9, 0.1]))\\
print(entropy([0.5, 0.25, 0.125, 0.125]))\\
rng = np.random.default\_rng(0)\\
H = [entropy(rng.dirichlet(np.ones(4)))\\
\hspace*{2em}\hspace*{2em}for \_ in range(10000)]\\
print(max(H))
\end{alltt}
        This prints $1.0$, $0.469$ and $1.75$, and the largest of the
        random entropies is $1.9998$: close to, but never above, 2 bits.
      \end{answer}}{\lines{6}}

  \end{qparts}
\end{question}
\newpage
% ***************************** question 2 **************************

\begin{question}{20}
  \begin{qparts}

    \qpart[5]{Bits $X$ and $Y$ have the joint distribution
      \begin{center}
        \begin{tabular}{c|cc}
          & $Y=0$ & $Y=1$ \\ \hline
          $X=0$ & 0.4 & 0.1 \\
          $X=1$ & 0.1 & 0.4
        \end{tabular}
      \end{center}
      Find $H_X$, $H_Y$, $H_{X,Y}$, $H_{X|Y}$ and the mutual information
      $I_{X;Y}$.
      \begin{answer}
        Both marginals are $(\frac12,\frac12)$, so $H_X = H_Y = 1$ bit.
        $H_{X,Y} = -2\times0.4\log_2 0.4 - 2\times0.1\log_2 0.1 \approx
        1.722$ bits.  By the chain rule $H_{X|Y} = H_{X,Y} - H_Y \approx
        0.722$ bits, and $I_{X;Y} = H_X - H_{X|Y} \approx 0.278$ bits.
      \end{answer}}{\lines{6}}

    \qpart[3]{Explain why this table describes sending equally likely
      bits down a binary symmetric channel.  What is the error rate $f$,
      and what is the channel capacity $1-H(f)$?
      \begin{answer}
        $\Prob{X=0} = \Prob{X=1} = \frac12$ and
        $\Prob{Y\neq X\mid X} = 0.1/0.5 = 0.2$ whichever bit is sent, so
        $f = 0.2$.  $C = 1 - H(0.2) = 1 - 0.722 = 0.278$ bits per bit
        sent: the mutual information of part (a), because equally likely
        inputs achieve the capacity.
      \end{answer}}{\lines{4}}

    \qpart[4]{Let $p$ be a fair coin and $q$ a coin with
      $\Prob{\text{head}} = 0.9$.  Find $\KL{p}{q}$ and $\KL{q}{p}$ in
      bits.  What does $\KL{p}{q}$ mean in terms of code lengths?
      \begin{answer}
        $\KL{p}{q} = \frac12\log_2\frac{0.5}{0.9} +
        \frac12\log_2\frac{0.5}{0.1} \approx 0.737$ bits and
        $\KL{q}{p} = 0.9\log_2\frac{0.9}{0.5} + 0.1\log_2\frac{0.1}{0.5}
        \approx 0.531$ bits.  They differ: the KL divergence is not
        symmetric (so it is not a distance).  $\KL{p}{q}$ is the expected
        number of extra bits per toss when we code a fair coin with the
        code that is optimal for the biased one.
      \end{answer}}{\lines{6}}

    \qpart[4]{For some input, the true label distribution over three
      classes is $p = (0.7, 0.2, 0.1)$ and a classifier predicts
      $q = (0.6, 0.3, 0.1)$.  Using natural logarithms, find the expected
      cross-entropy loss $\av[Y\sim p]{-\ln q(Y)}$, the entropy $H_p$
      and $\KL{p}{q}$, and check that the loss is $H_p + \KL{p}{q}$.
      \begin{answer}
        Loss $= -(0.7\ln0.6 + 0.2\ln0.3 + 0.1\ln0.1) \approx 0.829$ nats.
        $H_p = -(0.7\ln0.7 + 0.2\ln0.2 + 0.1\ln0.1) \approx 0.802$ nats.
        $\KL{p}{q} = 0.7\ln\frac{0.7}{0.6} + 0.2\ln\frac{0.2}{0.3} \approx
        0.027$ nats, and $0.802 + 0.027 = 0.829$.  Only the KL term depends
        on the classifier, so lowering the loss moves $q$ towards $p$; the
        loss can never go below $H_p$.
      \end{answer}}{\lines{6}}

    \qpart[4]{In Python, check part (d), then show numerically that the
      KL divergence is never negative by computing it for $10\,000$ random
      pairs of distributions on five outcomes.
      \begin{answer}
\begin{alltt}
p = np.array([0.7, 0.2, 0.1])\\
q = np.array([0.6, 0.3, 0.1])\\
ce = -(p * np.log(q)).sum()\\
h = -(p * np.log(p)).sum()\\
kl = (p * np.log(p / q)).sum()\\
print(ce, h, kl, h + kl)\\
rng = np.random.default\_rng(1)\\
P = rng.dirichlet(np.ones(5), 10000)\\
Q = rng.dirichlet(np.ones(5), 10000)\\
print((P * np.log(P / Q)).sum(1).min())
\end{alltt}
        This prints $0.829$, $0.802$, $0.027$ and $0.829$; the smallest of
        the random KL divergences is $0.005$, positive as it must be.
      \end{answer}}{\lines{6}}

  \end{qparts}
\end{question}
\newpage
% ***************************** question 3 **************************

\begin{question}{20}
  Maximum entropy with discrete variables needs only ordinary Lagrange
  multipliers.  We use natural logarithms here, so
  $H = -\sum_i p_i\ln p_i$.
  \begin{qparts}

    \qpart[4]{A dice has faces $i = 1,\ldots,6$ with probabilities
      $p_i$.  Maximise the entropy subject only to $\sum_i p_i = 1$, using
      the Lagrangian $\mathcal{L} = -\sum_i p_i\ln p_i +
      \lambda_0\bigl(\sum_i p_i - 1\bigr)$.  What distribution do you get?
      \begin{answer}
        $\partial\mathcal{L}/\partial p_i = -\ln p_i - 1 + \lambda_0 = 0$,
        so $p_i = \e{\lambda_0 - 1}$ is the same for every face, and
        normalising gives $p_i = \frac16$: with no information beyond
        normalisation, maximum entropy gives the uniform distribution.
      \end{answer}}{\lines{5}}

    \qpart[5]{Now we are told only that the average score is $\mu$:
      $\sum_i i\,p_i = \mu$.  Add a multiplier $\beta$ for this
      constraint and show that
      \begin{align*}
        p_i = \frac{\e{\beta i}}{Z(\beta)}, \qquad
        Z(\beta) = \sum_{j=1}^6 \e{\beta j}.
      \end{align*}
      \begin{answer}
        With $\mathcal{L} = -\sum_i p_i\ln p_i + \lambda_0(\sum_i p_i - 1)
        + \beta(\sum_i i\,p_i - \mu)$,
        \begin{align*}
          \frac{\partial\mathcal{L}}{\partial p_i} = -\ln p_i - 1 +
          \lambda_0 + \beta\,i = 0
          \quad\Rightarrow\quad p_i = \e{\lambda_0-1}\,\e{\beta i}.
        \end{align*}
        Normalising fixes $\e{\lambda_0-1} = 1/Z(\beta)$.  This is the
        softmax of the lecture with scores $s_i = i$ and temperature
        $T = 1/\beta$.
      \end{answer}}{\lines{8}}

    \qpart[3]{Which value of $\beta$ gives $\mu = 3.5$?  Should $\beta$
      be positive or negative if $\mu = 4.5$?  Explain.
      \begin{answer}
        $\beta = 0$ gives the uniform distribution, whose mean is $3.5$.
        For $\mu = 4.5$ the high faces must be more likely, so
        $\e{\beta i}$ must grow with $i$: $\beta > 0$.
      \end{answer}}{\lines{4}}

    \qpart[5]{In Python, find the $\beta$ that gives $\mu = 4.5$ (use
      \texttt{scipy.optimize.brentq} on $\mu(\beta) - 4.5$, or bisection),
      and print the probabilities and the entropy in bits.
      \begin{answer}
\begin{alltt}
from scipy.optimize import brentq\\
i = np.arange(1, 7)\\
def mean(b):\\
\hspace*{2em}w = np.exp(b * i)\\
\hspace*{2em}return (w * i).sum() / w.sum()\\
b = brentq(lambda b: mean(b) - 4.5, -5, 5)\\
p = np.exp(b * i) / np.exp(b * i).sum()\\
print(b, p, -(p * np.log2(p)).sum())
\end{alltt}
        $\beta \approx 0.371$ and $p \approx (0.054, 0.079, 0.114, 0.165,
        0.240, 0.347)$, with entropy $2.33$ bits against $\log_2 6 \approx
        2.585$ for a fair dice.
      \end{answer}}{\lines{6}}

    \qpart[3]{Why is this a sensible guess for the dice, and when might
      it be a poor one?
      \begin{answer}
        Of all distributions with mean $4.5$ it assumes the least: it is
        the most uncertain one consistent with what we know, and (from
        the counting argument in the lecture) by far the most likely
        outcome of many throws.  It is a poor guess if we know more, for
        instance that the dice is loaded only on one face.
      \end{answer}}{\lines{4}}

  \end{qparts}
\end{question}
\newpage

% ***************************** question 4 **************************

\begin{question}{25}
  This question is harder and optional.  We move a probability mass $p_i$
  at points $x_i$ onto a mass $q_j$ at points $y_j$, where moving a unit
  of mass from $x_i$ to $y_j$ costs $d_{ij}$.  A transport plan
  $\gamma_{ij}\geq0$ has row sums $\sum_j\gamma_{ij} = p_i$ and column
  sums $\sum_i\gamma_{ij} = q_j$, and the optimal transport cost is the
  smallest $\sum_{i,j} d_{ij}\gamma_{ij}$.
  \begin{qparts}

    \qpart[4]{In one dimension, with $d_{ij} = |x_i - y_j|$ and equal
      masses, the optimal plan matches the points in sorted order.  Find
      the cost of moving equal masses $\frac13$ at $x = (0, 1, 3)$ to
      $y = (2, 4, 5)$.  Is the optimal plan unique?
      \begin{answer}
        Matching in order, $W = \frac13(|0-2| + |1-4| + |3-5|) = \frac73
        \approx 2.333$.  It is not unique: matching $0\to4$, $1\to2$,
        $3\to5$ also costs $\frac13(4+1+2) = \frac73$.
      \end{answer}}{\lines{4}}

    \qpart[4]{The Sinkhorn algorithm sets $K_{ij} = \e{-d_{ij}/\epsilon}$,
      starts from $\bm{u} = \bm{1}$, and repeats
      $\bm{v}\leftarrow\bm{q}\oslash(\mat{K}^\tr\bm{u})$,
      $\bm{u}\leftarrow\bm{p}\oslash(\mat{K}\bm{v})$, where $\oslash$ is
      element-by-element division.  The plan is
      $\gamma_{ij} = u_i K_{ij} v_j$.  Show that straight after the
      $\bm{v}$ update the column sums of $\gamma$ are exactly $q_j$, and
      straight after the $\bm{u}$ update the row sums are exactly $p_i$.
      \begin{answer}
        $\sum_i u_i K_{ij} v_j = v_j\,(\mat{K}^\tr\bm{u})_j = q_j$, since
        $v_j = q_j/(\mat{K}^\tr\bm{u})_j$.  Likewise
        $\sum_j u_i K_{ij} v_j = u_i\,(\mat{K}\bm{v})_i = p_i$.  Each update
        fixes one set of sums and slightly upsets the other; repeating
        converges to a plan with both.
      \end{answer}}{\lines{6}}

    \qpart[7]{Implement Sinkhorn in PyTorch for part (a) (start with
      \texttt{torch.set\_default\_dtype(torch.float64)}: $\e{-d/\epsilon}$
      is tiny).  Run 1000 iterations for $\epsilon = 1$, $0.1$ and
      $0.01$, and print the cost $\sum_{ij}d_{ij}\gamma_{ij}$ and the row
      and column sums.  Look at the plan for $\epsilon = 0.01$.
      \begin{answer}
\begin{alltt}
import torch\\
torch.set\_default\_dtype(torch.float64)\\
x = torch.tensor([0.0, 1.0, 3.0])\\
y = torch.tensor([2.0, 4.0, 5.0])\\
p = torch.full((3,), 1/3)\\
q = torch.full((3,), 1/3)\\
d = (x[:, None] - y[None, :]).abs()\\
for eps in [1.0, 0.1, 0.01]:\\
\hspace*{2em}K = torch.exp(-d / eps)\\
\hspace*{2em}u = torch.ones(3)\\
\hspace*{2em}for \_ in range(1000):\\
\hspace*{4em}v = q / (K.T @ u)\\
\hspace*{4em}u = p / (K @ v)\\
\hspace*{2em}gamma = u[:, None] * K * v[None, :]\\
\hspace*{2em}print(eps, (gamma * d).sum().item())\\
\hspace*{2em}print(gamma.sum(1), gamma.sum(0))\\
print(gamma)
\end{alltt}
        The costs are $2.400$, $2.333$ and $2.333$: as $\epsilon\to0$ we
        recover $\frac73$.  The sums are all $\frac13$.  For
        $\epsilon = 0.01$ the plan is
        \begin{align*}
          \gamma \approx \begin{pmatrix} 1/6 & 1/12 & 1/12 \\ 1/6 & 1/12 &
          1/12 \\ 0 & 1/6 & 1/6 \end{pmatrix},
        \end{align*}
        not a matching.  Many plans have the optimal cost here, and the
        entropy term picks the most spread out of them.
      \end{answer}}{\lines{10}}

    \qpart[10]{Because every step is a PyTorch operation, the Sinkhorn
      cost can be differentiated.  Put your code in a function
      \texttt{sinkhorn\_cost(x, y, eps=0.1, iters=200)} for points in two
      dimensions, with equal masses and squared distances
      $d_{ij} = \len{\bm{x}_i - \bm{y}_j}^2$.  Take 30 target points
      $\bm{y}_j$ equally spaced on the circle of radius 1 about $(3,0)$,
      start 30 points $\bm{x}_i$ at random
      (\texttt{torch.randn(30, 2, requires\_grad=True)}), and run 300
      steps of \texttt{torch.optim.SGD([x], lr=0.5)} on the Sinkhorn cost.
      Where do the points end up?  Why squared distances?
      \begin{answer}
\begin{alltt}
def sinkhorn\_cost(x, y, eps=0.1, iters=200):\\
\hspace*{2em}d = ((x[:, None, :] - y[None, :, :])**2).sum(-1)\\
\hspace*{2em}K = torch.exp(-d / eps)\\
\hspace*{2em}p = torch.full((len(x),), 1/len(x))\\
\hspace*{2em}q = torch.full((len(y),), 1/len(y))\\
\hspace*{2em}u = torch.ones(len(x))\\
\hspace*{2em}for \_ in range(iters):\\
\hspace*{4em}v = q / (K.T @ u)\\
\hspace*{4em}u = p / (K @ v)\\
\hspace*{2em}return (u[:, None] * K * v[None, :] * d).sum()\\
torch.manual\_seed(0)\\
t = torch.linspace(0, 2*torch.pi, 31)[:-1]\\
y = torch.stack([3 + torch.cos(t), torch.sin(t)], 1)\\
x = torch.randn(30, 2, requires\_grad=True)\\
opt = torch.optim.SGD([x], lr=0.5)\\
for step in range(301):\\
\hspace*{2em}opt.zero\_grad()\\
\hspace*{2em}loss = sinkhorn\_cost(x, y)\\
\hspace*{2em}loss.backward()\\
\hspace*{2em}opt.step()\\
\hspace*{2em}if step \% 50 == 0:\\
\hspace*{4em}print(step, loss.item())
\end{alltt}
        The cost falls from $10.7$ to $0.051$, and the points end up on
        the circle (mean distance from $(3,0)$ $1.000$, spread $0.0004$),
        one near each target.  The small cost that remains comes from the
        entropy term.  The distance $\len{\bm{x}-\bm{y}}$ has no gradient
        when a point lands exactly on a target, while its square is smooth
        everywhere.  This is the idea behind training generative models
        with optimal transport: the points are the model's samples, and
        the gradient moves them towards the data.
      \end{answer}}{\lines{12}}

  \end{qparts}
\end{question}

\end{document}
